\documentclass[10pt,a4paper]{article}
\usepackage[utf8]{inputenc}
\usepackage[T1]{fontenc}
\usepackage[ngerman]{babel}
\usepackage[margin=1.7cm,top=1.5cm,bottom=1.7cm]{geometry}
\usepackage{titlesec}
\titleformat*{\section}{\normalfont\large\bfseries}
\titleformat*{\subsection}{\normalfont\normalsize\bfseries}
\titlespacing*{\section}{0pt}{8pt}{3pt}
\titlespacing*{\subsection}{0pt}{6pt}{2pt}
\usepackage{amsmath,amssymb}
\usepackage{listings}
\usepackage{xcolor}
\usepackage{tikz}
\usepackage{fancyhdr}
\usepackage{parskip}
\usepackage{needspace}

\newcommand{\mcbox}{\raisebox{-1pt}{\fbox{\rule{0pt}{7pt}\rule{7pt}{0pt}}}\hspace{6pt}}
\newcommand{\antwortlinien}[1]{%
  \par\vspace{2pt}%
  \newcount\zeilen \zeilen=#1
  \loop
    \noindent\rule{\linewidth}{0.4pt}\par\vspace{11pt}%
    \advance\zeilen by -1
  \ifnum\zeilen>0 \repeat
}
\lstset{language=Python,basicstyle=\ttfamily\footnotesize,keywordstyle=\bfseries,
  commentstyle=\itshape\color{gray},showstringspaces=false,frame=single,
  framerule=0.3pt,xleftmargin=8pt,framexleftmargin=6pt,aboveskip=6pt,belowskip=6pt}

\pagestyle{fancy}
\fancyhf{}
\fancyhead[L]{\small MDT5/2 -- Session 9: k-Means, Silhouette, PCA}
\fancyhead[R]{\small Übungsaufgaben zur Klausurvorbereitung}
\fancyfoot[C]{\small Seite \thepage}
\renewcommand{\headrulewidth}{0.4pt}

\begin{document}

\begin{center}
  {\Large\bfseries Übungsaufgaben Session 9}\\[4pt]
  {\large Kapitel 7a: k-Means, Ellbogen \& Silhouette, PCA, Kernel-PCA}\\[4pt]
  \small Stil der Übungssammlung MDT5/2 -- generierte Aufgaben, keine Originale
\end{center}

\vspace{4pt}

\section*{Aufgabe 1 -- Multiple Choice mit Begründung}
\emph{Markieren Sie jeweils die einzige richtige von drei Aussagen und begründen Sie Ihre Entscheidung.}

\subsection*{1.1 \ Für k-Means gilt:}
\mcbox k-Means findet Cluster beliebiger Form.

\mcbox Die Anzahl der Cluster muss vorgegeben werden; das Ergebnis sind tendenziell
kugelförmige Cluster um die Zentren.

\mcbox Eine Normalisierung der Merkmale ist für k-Means unwichtig.

Begründung:
\antwortlinien{2}

\subsection*{1.2 \ Für die Silhouette eines Punkts gilt:}
\mcbox Sie liegt im Intervall $[-1; 1]$; Werte nahe 1 bedeuten eine gute Zuordnung
zum eigenen Cluster.

\mcbox Sie liegt im Intervall $[0; \infty)$; hohe Werte deuten auf Ausreißer.

\mcbox Sie vergleicht den Abstand zum eigenen Cluster-Zentrum mit dem Radius des Clusters.

Begründung:
\antwortlinien{2}

\subsection*{1.3 \ Für die PCA gilt:}
\mcbox Die Hauptachsen sind orthogonal und zeigen in die Richtungen der größten Varianz
der Daten.

\mcbox Die PCA maximiert die Kovarianzen zwischen den transformierten Merkmalen.

\mcbox Die PCA funktioniert unabhängig davon, ob die Daten normalisiert sind.

Begründung:
\antwortlinien{2}

\section*{Aufgabe 2 -- k-Means verstehen}

a) Geben Sie die Schritte des k-Means-Algorithmus in der richtigen Reihenfolge an.
\antwortlinien{4}

b) Warum ist die rein zufällige Initialisierung der Cluster-Zentren problematisch,
und wie funktioniert k-means++ in Grundzügen?
\antwortlinien{4}

c) Beschreiben Sie, wie Sie mit einem trainierten k-Means-Clustering (i) Novelty Detection
und (ii) Outlier Detection in einem unbekannten Datensatz betreiben können.
\antwortlinien{4}

\section*{Aufgabe 3 -- Anzahl der Cluster bestimmen}

a) Für $k = 1, \dots, 6$ wurden folgende Werte der Summe der quadrierten Abstände
(\texttt{inertia\_}) gemessen:

\begin{center}
\begin{tabular}{c|cccccc}
$k$ & 1 & 2 & 3 & 4 & 5 & 6\\
\hline
inertia & 1000 & 400 & 150 & 130 & 120 & 115
\end{tabular}
\end{center}

Skizzieren Sie den Ellbogen-Plot und bestimmen Sie begründet die optimale Clusteranzahl.

\vspace{4cm}

b) Die Silhouette eines Punkts ist $S(\boldsymbol{p}) = \frac{b-a}{\max(a,b)}$.
Berechnen Sie $S$ für einen Punkt mit $a = 1$ (mittlerer Abstand zu Punkten im eigenen
Cluster) und $b = 3$ (mittlerer Abstand zu Punkten im nächsten Cluster) und
interpretieren Sie das Ergebnis.
\antwortlinien{2}

c) Berechnen und interpretieren Sie $S$ für einen Punkt mit $a = 2$ und $b = 1$.
\antwortlinien{2}

d) Woran erkennen Sie in einem Silhouetten-Plot (Silhouetten aller Punkte, gruppiert
nach Clustern, plus gestrichelte Linie für den Gesamtkoeffizienten) eine schlechte
Clusteranzahl? Nennen Sie zwei Anzeichen.
\antwortlinien{3}

e) Warum ist die Ellbogen-Methode \glqq intuitiv, aber subjektiv\grqq, und warum werden
in der Praxis häufig mehrere Verfahren kombiniert?
\antwortlinien{2}

\section*{Aufgabe 4 -- PCA anwenden}

a) Geben Sie die vier Anwendungs-Schritte der PCA aus der Vorlesung an
(Zentrierung, Datenmatrix, \dots).
\antwortlinien{4}

b) Eine PCA liefert die Eigenwerte $5$; $3$; $1{,}5$; $0{,}4$; $0{,}1$.
Wie viele Dimensionen $k$ müssen erhalten bleiben, damit mindestens 95\,\% der Varianz
erhalten bleiben? Rechnen Sie vor.
\antwortlinien{3}

c) Warum müssen die Daten vor der PCA normalisiert werden?
\antwortlinien{2}

d) Wie wird PCA zur Novelty Detection eingesetzt? Beschreiben Sie Training und
Score-Berechnung für neue Daten.
\antwortlinien{4}

e) Nennen Sie zwei Annahmen der PCA, bei deren Verletzung die Projektion keine gute
Trennung mehr ermöglicht -- und welche Erweiterung aus der Vorlesung hilft dann?
\antwortlinien{3}

\section*{Aufgabe 5 -- Clustering-Verfahren zu Plot beurteilen}

Die Daten enthalten drei bekannte Cluster: zwei kompakte runde Cluster und einen
langgestreckten, gebogenen Cluster (\glqq Banane\grqq), der sich um einen der runden legt.

\begin{center}
\begin{tikzpicture}[scale=0.7]
\draw[->] (-0.5,0) -- (10.5,0) node[right] {$x_1$};
\draw[->] (0,-0.5) -- (0,6.5) node[above] {$x_2$};
% rundes Cluster 1
\foreach \p in {(2,2),(2.3,2.2),(1.8,2.4),(2.5,1.9),(2.1,1.7),(1.7,2.0),(2.4,2.5),(2.0,2.8)}
  \fill \p circle (2pt);
% rundes Cluster 2 rechts
\foreach \p in {(8.5,1.5),(8.8,1.7),(8.3,1.9),(9.0,1.4),(8.6,1.2),(8.2,1.5),(8.9,2.0),(8.5,2.2)}
  \fill \p circle (2pt);
% Banane um Cluster 1
\foreach \a in {200,220,...,340} {
  \fill ({2+2.4*cos(\a)},{4.2+2.4*sin(\a)*(-1)}) circle (2pt);
  \fill ({2+2.7*cos(\a)},{4.2-2.7*sin(\a)}) circle (2pt);
}
\end{tikzpicture}
\end{center}

Begründen Sie, welche Erfolgsaussichten (i) k-Means und (ii) ein GMM haben, die drei
Cluster zu identifizieren. Was wäre das typische Fehlverhalten von k-Means bei der Banane?

\antwortlinien{6}

\vspace{10pt}
\begin{center}
\footnotesize\itshape
Nach Bearbeitung: Antworten an Claude schicken (Foto genügt) -- Korrektur mit Begründungs-Check.
\end{center}

\end{document}
